\chapter{Content and Assets}
\label{ch:content-and-assets}

Real XNA's content model and CNA's are different in one foundational way, and that
difference is a deliberate design choice this book has already flagged twice in passing
(Chapter~\ref{ch:what-is-cna}, Chapter~\ref{ch:first-game}). This chapter is where it gets a
full explanation.

\section{No \texttt{.xnb} content pipeline, by design}

Real XNA never ships raw asset files inside a game. A separate, offline content-pipeline
build step compiles every texture, model, font, and effect into a proprietary binary format
--- \texttt{.xnb} --- which the runtime's \cnaclass{ContentManager} loads via
\texttt{Load<T>()}. FNA reimplements a real \texttt{.xnb} reader, precisely so unmodified,
historical XNA content keeps working unchanged. CNA takes a different, explicitly-stated
path instead: it loads raw asset files directly (a \texttt{.png}, a \texttt{.json}
descriptor, a \texttt{.dds}) rather than a compiled binary blob. This is named in the
project's own README as a deliberate design choice, not an unfinished feature --- the same
status as the \texttt{.fx} bytecode gap from Chapter~\ref{ch:shader-fx-gap}, but a separate
gap with a separate story.

\section{Two loaders exist, and it matters which one your asset goes through}

Chapter~\ref{ch:models-meshes} already covered this in detail for \cnaclass{Model}
specifically, and the pattern generalizes: for several asset types, CNA actually has
\emph{two} independent loading paths, and \texttt{ContentManager} tries the more capable one
first.

For \cnaclass{Model}, a real, wire-compatible \texttt{.xnb} \texttt{ModelReader} exists ---
genuinely compatible with real XNA/MonoGame/FNA-compiled model assets, reconstructing the
full bone hierarchy, per-mesh parent-bone assignment, a real bounding sphere, and correct
shared-resource deduplication across meshes. This is not a partial or aspirational reader;
where it applies, none of this chapter's caveats apply. Only when that real reader does not
recognize an asset does \texttt{ContentManager} fall back to an older, loose-file
\texttt{.model.json} loader, and it is this second loader that carries every one of the real
gaps catalogued in Chapter~\ref{ch:models-meshes}: a flattened single-bone hierarchy, no
per-mesh parent-bone assignment, a permanently degenerate bounding sphere, and no shared
resource graph across meshes sharing geometry.

This same two-tier shape --- a genuine, more-capable loader tried first, an older, simpler
loader as fallback --- recurs elsewhere in the content system, and the practical implication
for a porting engineer is the same each time: \emph{which} loader your asset actually goes
through matters more than what the target class's API surface promises in the abstract.

\section{JSON descriptors: the shape of the fallback path}

Where CNA does not have a real binary-compatible reader for an asset type, the fallback
convention is a hand-authored JSON descriptor sitting alongside the raw asset file ---
readable, editable by hand, and versioned in an ordinary source tree rather than requiring a
separate content-pipeline build step to regenerate. This is a genuinely different authoring
workflow from real XNA's content pipeline, and Volume~II's migration chapter returns to what
it means in practice: porting an existing XNA/FNA game means translating its compiled
\texttt{.xnb} assets into this raw-file-plus-JSON shape (or, for \cnaclass{Model}
specifically, discovering that the real \texttt{.xnb} reader already accepts the original
asset unmodified), rather than assuming every asset ports as a literal file copy.

\section{A worked example: content in Chapter~\ref{ch:first-game}'s smallest program}

The one-line \texttt{Texture2D::FromStream}-style construction in
Chapter~\ref{ch:first-game}'s minimal game --- \texttt{Texture2D("assets/logo.png",
getGraphicsDeviceProperty())} --- is possible precisely because of this chapter's design
choice: \texttt{"assets/logo.png"} is a literal, directly loadable file path, not a compiled
content-pipeline output identifier the way a real XNA \texttt{Content.Load<Texture2D>%
("logo")} call would be. This is simpler for a small, from-scratch CNA project, and it is
exactly the substitution a porting engineer has to reverse when bringing an existing XNA
asset pipeline's output over --- covered in full in Volume~II.

\section{Where this leaves the rest of Part~III}

Every class in the chapters ahead --- \cnaclass{Texture2D}, \cnaclass{Model}, the stock
effects --- is describable purely in terms of its runtime API, independent of how an
instance of it actually got constructed. This chapter's distinction matters specifically at
the boundary between ``I have an asset file'' and ``I have a live CNA object,'' which is
exactly the boundary Volume~II's migration guide and \texttt{cna-samples}/\texttt{cna-examples}
chapter come back to when the asset in question was originally authored for real XNA rather
than for CNA from scratch.
